\chapter{Literature Survey}

Several researchers have contributed significant work in the field of image super-resolution. Various deep learning-based models such as CNNs and GANs have been proposed to enhance image quality by learning high-frequency features. The following studies represent some of the prominent developments in this area.
\vspace{0.5cm}

\textbf{CNN-based Super-Resolution:}

Dong et al. (2014) \cite{a2} introduced SRCNN, the first deep learning-based SR model using convolutional neural networks (CNNs). SRCNN demonstrated significant performance improvements over previous methods in terms of PSNR and SSIM. It was later followed by deeper architectures such as VDSR (Kim et al., 2016) and DRCN, which leveraged residual learning to further enhance results. \cite{a3}

\vspace{0.3cm}

\textbf{GAN-based Super-Resolution:}


\textbf{Lai et al. (2017)} \cite{a10}Proposed a Laplacian pyramid network structure for progressive image reconstruction with improved accuracy and speed.(LapSRN)

\textbf{Ledig et al. (2017)} \cite{a4} proposed the SRGAN model, which marked a major milestone in perceptual super-resolution. SRGAN combined an adversarial loss with perceptual loss derived from VGG networks, resulting in photo-realistic high-resolution images.

\textbf{Lim et al. (2017)} \cite{a9}Removed unnecessary modules from traditional residual networks, allowing for deeper and more efficient SR models.(EDSR) 

\textbf{Wang et al. (2018)} \cite{a5}Enhanced the SRGAN framework by introducing Residual-in-Residual Dense Blocks (RRDB) and a relativistic discriminator.(ESRGAN) 

\textbf{Zhang et al. (2018)} \cite{a8}Introduced a channel attention mechanism to adaptively rescale channel-wise features for improved high-frequency detail restoration.(RCAN) 

\textbf{Liang et al. (2021)} \cite{a7}Utilized Swin Transformers for SR, achieving state-of-the-art results on multiple benchmarks.(SwinIR)


\textbf{Wang et al. (2021)} \cite{a6}Aimed at real-world image SR by handling unknown degradations using a generalized degradation model.(Real-ESRGAN)

\textbf{Chai et al. (2022)} \cite{a11} Proposed TPE-GAN, a GAN-based method for image encryption that preserves thumbnails while ensuring security through key-based access. (TPE-GAN)

\textbf{Fang et al. (2022)} \cite{a12} Introduced Atten-GAN for pedestrian trajectory prediction by incorporating attention mechanisms into GANs to improve contextual understanding. (Atten-GAN)